\documentclass[12pt,letterpaper]{article}

% ==========================================================
% PLANTILLA LATEX / ICONTEC
% Curso Metodología de la Investigación - UPTC
% Generada a partir del Instrumento de trabajo (Lección 1):
% "Selección de la idea de investigación" - únicamente la
% información propia del equipo (no el ejemplo resuelto).
% ==========================================================

\usepackage[top=3cm,left=4cm,bottom=3cm,right=2cm]{geometry}
\usepackage[utf8]{inputenc}
\usepackage[T1]{fontenc}
\usepackage{mathptmx} % tipografía tipo Times, similar al PDF de referencia
\usepackage{setspace}
\usepackage{titlesec}
\usepackage{tabularx}
\usepackage{longtable}
\usepackage{booktabs}
\usepackage{array}
\usepackage[table]{xcolor}
\usepackage[most]{tcolorbox}
\usepackage{enumitem}
\usepackage{hyperref}
\hypersetup{colorlinks=false, hidelinks}
\usepackage{xurl}

\onehalfspacing

% ---------- Colores (tono ámbar del badge "EN CONSTRUCCIÓN") ----------
\definecolor{badgebg}{HTML}{FDF6E3}
\definecolor{badgeborder}{HTML}{C9A227}
\definecolor{badgetext}{HTML}{8A6D1D}
\definecolor{greensignal}{HTML}{2E7D32}

% ---------- Formato de títulos de sección (centrado, versalitas, como el PDF) ----------
\titleformat{\section}{\centering\bfseries\large}{\thesection.}{0.5em}{\MakeUppercase}
\titleformat{\subsection}{\bfseries\normalsize}{\thesubsection}{0.5em}{}
\titlespacing*{\section}{0pt}{2.5em}{1.5em}

% ---------- Comando para el "badge" EN CONSTRUCCIÓN / estado ----------
\newcommand{\estado}[1]{%
  \begin{center}
  \tcbox[colback=badgebg,colframe=badgeborder,boxrule=0.5pt,
    arc=6pt,left=6pt,right=6pt,top=3pt,bottom=3pt,
    fontupper=\small\bfseries\color{badgetext}]{\MakeUppercase{#1}}
  \end{center}
}

% ---------- Comando para la nota itálica que acompaña cada estado ----------
\newcommand{\notafase}[1]{\begin{center}\itshape\small #1\end{center}}

\newcolumntype{L}[1]{>{\raggedright\arraybackslash}p{#1}}

\begin{document}

% ================= PORTADA =================
\begin{titlepage}
\begin{center}
\vspace*{1.5cm}
{\large UNIVERSIDAD PEDAGÓGICA Y TECNOLÓGICA DE COLOMBIA}\\[0.3cm]
{\large FACULTAD DE INGENIERÍA}\\[0.3cm]
{\large INGENIERÍA DE SISTEMAS Y COMPUTACIÓN}

\vspace{3cm}

{\Large\bfseries RELACIÓN ENTRE LAS HORAS DE EXPOSICIÓN A PANTALLAS,\\
LOS HÁBITOS POSTURALES Y LOS SÍNTOMAS FÍSICOS EN\\
ESTUDIANTES DE INGENIERÍA DE SISTEMAS DE LA UPTC}

\vspace{0.8cm}
{\itshape\small Título provisional --- se ajusta al formular la propuesta (Fase 3)}

\vspace{3cm}

{\bfseries AUTORES}\\[0.6cm]
Gabriel Alejandro Gamba Leguizamón\\
José Fernando Barreto Franco\\
Daniel Felipe Molina Barajas

\vfill

Documento elaborado en el marco del curso Metodología de la Investigación\\[0.3cm]
Tunja, Colombia · 2026

\end{center}
\end{titlepage}

% ================= CONTENIDO (tabla de estado por sección) =================
\section*{\centering CONTENIDO}
\begin{center}
\begin{tabularx}{\textwidth}{L{10cm} X}
\toprule
\textbf{Sección} & \textbf{Estado}\\
\midrule
Índice de tablas & incluido\\
Índice de figuras & incluido\\
1. Planteamiento del problema & parcial\\
\quad 1.1 Antecedentes & en construcción\\
\quad 1.2 Descripción de la situación problemática & redactado\\
\quad 1.3 Formulación del problema & en construcción\\
\quad 1.4 Consecuencias de no investigar & en construcción\\
\quad 1.5 Delimitación & en construcción\\
2. Justificación & en construcción\\
3. Objetivos & en construcción\\
4. Alcance de la investigación & en construcción\\
5. Formulación de hipótesis & en construcción\\
6. Variables & en construcción\\
7. Marco referencial & en construcción\\
8. Diseño metodológico & en construcción\\
9. Resultados y discusión & en construcción\\
10. Conclusiones y recomendaciones & en construcción\\
Referencias & redactado (parcial)\\
Anexos & incluido (parcial)\\
\bottomrule
\end{tabularx}
\end{center}

\newpage
% ================= ÍNDICE DE TABLAS =================
\section*{\centering ÍNDICE DE TABLAS}
Tabla 1. Relación entre síntomas y causas\dotfill 1\\
Tabla 2. Relación entre causas y síntomas\dotfill 1

\vspace{1cm}
\notafase{Esta lista crece en fases posteriores.}

\newpage
% ================= ÍNDICE DE FIGURAS =================
\section*{\centering ÍNDICE DE FIGURAS}
\notafase{Aún no hay figuras incorporadas al cuerpo del documento. Se listan aquí a medida que cada fase incorpore figuras propias (diagramas, gráficos, esquemas, etc.).}

\newpage
% ================= 1. PLANTEAMIENTO DEL PROBLEMA =================
\section{Planteamiento del problema}

\subsection{Antecedentes}
\estado{En construcción}
\notafase{Se redacta en la Fase 2 del curso.}

\subsection{Descripción de la situación problemática}

La problemática de esta investigación se relaciona con las condiciones en que los
estudiantes de Ingeniería de Sistemas desarrollan sus actividades académicas,
especialmente aquellas que requieren un uso prolongado del computador y otros
dispositivos electrónicos. Las extensas jornadas dedicadas al estudio, la
programación, la consulta de información y el desarrollo de proyectos pueden generar
molestias físicas y visuales, como fatiga ocular, sequedad o irritación de los ojos,
dolores o molestias en el cuello, hombros, espalda, brazos y muñecas, además de
cansancio físico y mental, particularmente cuando estas actividades se realizan
acompañadas de posturas inadecuadas, permanencia prolongada en una misma posición y
pausas insuficientes [1], [2]. En población universitaria se ha reportado una elevada
frecuencia de estas afectaciones; en particular, el síndrome visual informático
constituye una problemática frecuente entre estudiantes expuestos de manera
prolongada a dispositivos electrónicos [2], mientras que los trastornos
musculoesqueléticos presentan una alta frecuencia entre universitarios, con
afectaciones principalmente en regiones como el cuello, la zona lumbar y la espalda
[1]. Esta situación puede verse favorecida por la exposición continua a pantallas,
debido al esfuerzo visual sostenido y a la ausencia de periodos adecuados de
descanso; por la adopción de posturas inadecuadas, como mantener el cuello inclinado
hacia adelante, la espalda encorvada, los hombros elevados o las muñecas en
posiciones poco favorables; y por la permanencia durante largos periodos en una
misma posición. Asimismo, la falta de pausas activas y descansos visuales puede
reducir las oportunidades de recuperación física durante las jornadas académicas,
mientras que unas condiciones ergonómicas inadecuadas del puesto de estudio, como
mobiliario que no se ajusta correctamente, computadores portátiles ubicados a una
altura inconveniente o dispositivos de entrada utilizados en posiciones poco
favorables, pueden contribuir a la aparición de molestias musculoesqueléticas [1]. A
estas condiciones se suman factores relacionados con el entorno de estudio, como una
iluminación inadecuada, reflejos o una configuración poco apropiada de la pantalla,
que pueden incrementar el esfuerzo visual durante sesiones prolongadas frente a
dispositivos electrónicos [2]. Como consecuencia de estas condiciones, los
estudiantes pueden experimentar fatiga visual, manifestada mediante sensación de ojos
cansados, secos, irritados o visión borrosa; molestias o dolor en cuello y hombros
asociados con posturas estáticas y prolongadas; molestias en espalda, brazos, manos y
muñecas relacionadas con la permanencia prolongada frente al computador y con
posiciones poco adecuadas durante el uso del teclado y el mouse; así como cansancio
físico y mental, que puede afectar el bienestar y la concentración durante las
actividades académicas [1], [3]. De igual manera, la exposición prolongada a
dispositivos electrónicos puede relacionarse con dolores de cabeza y otras molestias
asociadas al uso intensivo de pantallas [3]. Aunque estas manifestaciones pueden
presentarse con frecuencia durante la formación universitaria, algunos estudiantes
pueden llegar a normalizarlas como parte de las exigencias de su carrera y no
reconocer oportunamente la importancia de adoptar medidas preventivas. En este
contexto, la realización de pausas activas y cambios periódicos de posición
constituye una estrategia que puede favorecer la recuperación durante las jornadas
académicas y contribuir al bienestar físico de los estudiantes [4]. Por esta razón,
resulta pertinente analizar la relación existente entre las condiciones de estudio,
los hábitos frente al computador y las molestias físicas y visuales que pueden
experimentar los estudiantes de Ingeniería de Sistemas, con el propósito de
identificar los principales factores asociados y contribuir al planteamiento de
medidas preventivas que favorezcan condiciones más adecuadas para el desarrollo de
sus actividades académicas.

\bigskip
\begin{center}
\begin{longtable}{c L{7cm} L{5cm}}
\caption{Relación entre síntomas y causas} \\
\toprule
\textbf{Ord.} & \textbf{Síntoma} & \textbf{Causas relacionadas} \\
\midrule
\endfirsthead
1 & Fatiga visual después de jornadas prolongadas frente a pantallas & 1, 2, 3, 4 \\
2 & Sensación de ojos secos, cansados o irritados & 1, 2, 3 \\
3 & Dolor o molestia en cuello y hombros & 1, 3, 4, 5 \\
4 & Dolor o molestia en espalda & 1, 3, 4, 5 \\
5 & Molestias en muñecas, brazos o manos & 1, 3, 4, 5, 6 \\
6 & Dolores de cabeza y disminución de la comodidad durante el estudio & 1, 2, 3 \\
\bottomrule
\end{longtable}
\end{center}

\bigskip
\begin{center}
\begin{longtable}{c L{7cm} L{5cm}}
\caption{Relación entre causas y síntomas} \\
\toprule
\textbf{Ord.} & \textbf{Causa} & \textbf{Síntomas relacionados} \\
\midrule
\endfirsthead
1 & Exposición prolongada a pantallas durante las jornadas de estudio & 1, 2, 3, 4, 5 \\
2 & Falta de pausas o descansos visuales durante el estudio & 1, 2 \\
3 & Permanecer durante largos periodos en una misma posición & 1, 2, 3, 4, 5 \\
4 & Adopción de posturas inadecuadas durante el uso del computador & 3, 4, 5, 6 \\
5 & Puesto de estudio con condiciones ergonómicas inadecuadas & 3, 4, 5, 6 \\
6 & Falta de pausas activas y cambios de posición & 3, 4, 5, 6 \\
\bottomrule
\end{longtable}
\end{center}

\subsection{Formulación del problema}
\estado{En construcción}
\notafase{Se redacta en la Fase 2 del curso.}

\subsection{Consecuencias de no investigar}
\estado{En construcción}
\notafase{Se redacta en la Fase 2 del curso.}

\subsection{Delimitación}
\estado{En construcción}
\notafase{Se redacta en la Fase 2 del curso.}

\newpage
% ================= 2. JUSTIFICACIÓN =================
\section{Justificación}
\estado{En construcción}
\notafase{Se redacta en la Fase 3 del curso.}

\newpage
% ================= 3. OBJETIVOS =================
\section{Objetivos}
\estado{En construcción}
\notafase{Se redacta en la Fase 3 del curso.}

\newpage
% ================= 4. ALCANCE =================
\section{Alcance de la investigación}
\estado{En construcción}
\notafase{Se redacta en la Fase 3 del curso.}

\newpage
% ================= 5. HIPÓTESIS =================
\section{Formulación de hipótesis}
\estado{En construcción}
\notafase{Se redacta en la Fase 3 del curso.}

\newpage
% ================= 6. VARIABLES =================
\section{Variables}
\estado{En construcción}
\notafase{Se redacta en la Fase 3 del curso.}

\newpage
% ================= 7. MARCO REFERENCIAL =================
\section{Marco referencial}
\estado{En construcción}
\notafase{Se redacta en la Fase 4 del curso.}

\newpage
% ================= 8. DISEÑO METODOLÓGICO =================
\section{Diseño metodológico}
\estado{En construcción}
\notafase{Se redacta en la Fase 5 del curso.}

\newpage
% ================= 9. RESULTADOS =================
\section{Resultados y discusión}
\estado{En construcción}
\notafase{Se redacta en la Fase 6 del curso.}

\newpage
% ================= 10. CONCLUSIONES =================
\section{Conclusiones y recomendaciones}
\estado{En construcción}
\notafase{Se redacta en la Fase 7 del curso.}

\newpage
% ================= REFERENCIAS =================
\section*{\centering REFERENCIAS}

\begin{enumerate}[leftmargin=1.8em,label=\lbrack\arabic*\rbrack,itemsep=0.6em]
\item F. Ccami-Bernal, P. Urday-Ramos, F. Zela-Coila, J. K. Cáceres-Ruiz y V. Cabrera-Caso, ``Trastornos musculoesqueléticos y prácticas ergonómicas en universitarios peruanos durante la pandemia de la COVID-19,'' \emph{Revista Cubana de Ortopedia y Traumatología}, vol. 37, no. 3, 2023. [En línea]. Disponible en: \url{https://revortopedia.sld.cu/index.php/revortopedia/article/view/608}

\item D. Fernández-Villacorta, A. N. Soriano-Moreno, T. Gálvez-Olórtegui, N. Agüero-Santivañez, D. R. Soriano-Moreno y V. A. Benites-Zapata, ``Síndrome visual informático en estudiantes universitarios de posgrado de una universidad privada de Lima, Perú,'' \emph{Archivos de la Sociedad Española de Oftalmología}, vol. 96, no. 10, 2021. [En línea]. Disponible en: \url{https://www.elsevier.es/es-revista-archivos-sociedad-espanola-oftalmologia-296-articulo-sindrome-visual-informatico-estudiantes-universitarios-S0365669121000058}

\item R. M. Meneses Castañeda, S. L. Ramos Rodríguez, C. D. C. Molfino Jaramillo, E. L. Sánchez Miraval, D. F. Stein Montoros y L. G. Chávez Rodríguez, ``Frecuencia del síndrome visual informático en estudiantes de sexto año de medicina durante la educación virtual por COVID-19,'' \emph{Revista de la Facultad de Medicina Humana}, 2021.

\item Zhunio Suin et al., ``Pausas activas y su importancia en el bienestar de estudiantes universitarios,'' 2024.
\end{enumerate}

\notafase{Lista numerada por orden de primera aparición en el texto; crece en fases posteriores.}

\newpage
% ================= ANEXOS =================
\section*{\centering ANEXOS}

\subsection*{Anexo A. Instrumento de selección de idea de investigación (Lección 1)}
\notafase{Información diligenciada por el equipo en el instrumento de la Lección 1.}

\bigskip
\noindent\textbf{Datos del equipo}
\begin{center}
\begin{tabularx}{\textwidth}{L{5cm} X}
\toprule
\textbf{Campo} & \textbf{Respuesta}\\
\midrule
Integrante 1 & Gabriel Alejandro Gamba Leguizamón\\
Integrante 2 & José Fernando Barreto Franco\\
Integrante 3 & Daniel Felipe Molina Barajas\\
Fecha & 23 de agosto de 2026\\
Tema de interés general & Salud Digital, Ergonomía y Bienestar Físico en Estudiantes de Ingeniería de Sistemas\\
\bottomrule
\end{tabularx}
\end{center}

\bigskip
\noindent\textbf{Idea de investigación tentativa}

\noindent Quiero saber cuál es la relación entre el número de horas diarias
dedicadas a la exposición a pantallas, las posturas adoptadas durante las jornadas
de estudio y la frecuencia de síntomas físicos (fatiga visual y molestias
osteomusculares) en los estudiantes de Ingeniería de Sistemas de la UPTC, porque
eso importaría para proponer una guía práctica de ergonomía y pausas activas
adaptada a las rutinas de los futuros profesionales de software.

\bigskip
\noindent\textbf{Interés y pertinencia}
\begin{center}
\begin{tabularx}{\textwidth}{L{5.5cm} X}
\toprule
\textbf{Pregunta} & \textbf{Respuesta del equipo}\\
\midrule
¿Cuál es su interés personal en este tema? & Como estudiantes de Ingeniería de Sistemas, pasamos largas jornadas frente al computador y frecuentemente experimentamos fatiga visual o dolores de espalda y cuello, por lo que queremos entender cómo nuestras rutinas afectan la salud física.\\
¿Cuál es su interés académico o profesional? & El tema se relaciona directamente con la salud ocupacional e higiene postural en el ámbito de la ingeniería de software, permitiéndonos adquirir conocimientos sobre la prevención de riesgos físicos derivados del trabajo sedentario y continuo con computadores.\\
¿Es un tema prioritario o significativo para la disciplina? & Sí. La extensión de jornadas de programación y clases virtuales aumenta la incidencia de síndromes como el Síndrome Visual Informático y trastornos musculoesqueléticos en estudiantes e ingenieros, impactando su rendimiento y calidad de vida.\\
¿A quién afecta el problema que identificaron? & Principalmente a los estudiantes de la Escuela de Ingeniería de Sistemas, docentes, y profesionales del área que pasan más de 6 horas diarias expuestos a pantallas.\\
\bottomrule
\end{tabularx}
\end{center}

\bigskip
\noindent\textbf{Originalidad}
\begin{center}
\begin{tabularx}{\textwidth}{L{5.5cm} X}
\toprule
\textbf{Pregunta} & \textbf{Respuesta del equipo}\\
\midrule
¿Ya se ha investigado esto antes? & Existen estudios generales sobre ergonomía y uso de pantallas en trabajadores de oficina, pero pocos enfocados específicamente en la población universitaria de Ingeniería de Sistemas de nuestro contexto regional.\\
Si ya existe, ¿qué preguntas quedan sin responder? & Cómo se relacionan específicamente el número acumulado de horas continuas de programación con la aparición temprana de molestias osteomusculares y visuales en estudiantes jóvenes.\\
¿Qué aportarían ustedes que no exista ya? & Datos primarios y actualizados del contexto universitario local (UPTC) junto con un análisis correlacional sencillo entre hábitos posturales, horas de pantalla y síntomas físicos reportados.\\
\bottomrule
\end{tabularx}
\end{center}

\bigskip
\noindent\textbf{Viabilidad}
\begin{center}
\begin{tabularx}{\textwidth}{L{5.5cm} X}
\toprule
\textbf{Pregunta} & \textbf{Respuesta del equipo}\\
\midrule
¿Qué tan complejo es el tema para desarrollarlo en 16 semanas? & Muy viable y acotado. Consiste en diseñar y aplicar una encuesta estructurada (usando escalas validadas como la escala de usabilidad o dolor sintomático) y analizar los datos correlacionales.\\
¿Tienen acceso a los datos, herramientas o participantes que necesitarían? & Sí. Tenemos acceso directo a nuestros compañeros de carrera como participantes del estudio y utilizaremos herramientas gratuitas como Google Forms para la recolección de datos y programas como Excel para el análisis de frecuencia.\\
¿Tienen (o pueden adquirir en el curso) los conocimientos técnicos necesarios? & Sí. El diseño de encuestas y el análisis estadístico descriptivo/correlacional básico se abordan en el marco del curso de Metodología de la Investigación.\\
¿El tiempo disponible (4 h de clase + 4 h autónomas por semana) alcanza para esta idea? & Sí. El tiempo es suficiente para la revisión de antecedentes, diseño del cuestionario, recolección de respuestas en 2 semanas, análisis estadístico y redacción del informe final.\\
\bottomrule
\end{tabularx}
\end{center}

\bigskip
\noindent\textbf{Semáforo de la idea}
\begin{center}
\begin{tabularx}{\textwidth}{L{2.5cm} L{6cm} X}
\toprule
\textbf{Señal} & \textbf{Significa} & \textbf{¿Aplica a la idea?}\\
\midrule
\textcolor{greensignal}{\textbf{Verde}} & Cumple los tres criterios sin ajustes mayores. & \textbf{Sí} --- puede avanzar a la Lección 2 (antecedentes).\\
Amarillo & Cumple en general, pero un criterio necesita ajustarse. & No\\
Rojo & Falla claramente en al menos un criterio. & No\\
\bottomrule
\end{tabularx}
\end{center}

\end{document}
